\section{When the share of spam changes}
\label{sec:shift}
A filter trained where $39\%$ of emails are spam may be used where only $5\%$ are. Its probabilities
then overstate spam, and the thresholds of Section~\ref{sec:rule} act on the wrong numbers.

\subsection{Correcting the probabilities}
\label{sec:t5}
\textbf{Code:} \code{prior\_shift.py}, \code{experiments/t5\_prior\_correction.py}.

If spam still looks the same and only its share changes from $\pio$ to $\pion$, every email's
log-odds should move by the same amount, so \code{prior\_shift.py} corrects the model's output as
\[
  \post'=\sigm\!\big(\logit\post+\lshift\big),\qquad
  \lshift=\log\frac{\pion}{1-\pion}-\log\frac{\pio}{1-\pio}.
\]
For example, a model trained at $\pio=0.5$ and deployed at $\pion=0.1$ turns $\post=0.98$ into
$\post'\approx0.845$, so that email is inspected instead of filed as spam. Equivalently, the
thresholds can be moved instead of the probabilities (Figure~\ref{fig:t5}).

\begin{figure}[!htbp]
\centering
\includegraphics[width=0.92\linewidth]{t5_prior_correction.pdf}
\caption{Prior-shift correction for a model trained at $\pio=0.5$. Left: corrected $\post'$ against
$\post$. Right: the equivalent thresholds on the uncorrected $\post$ for every deployment rate $\pion$;
the shaded band is the inspect region.}
\label{fig:t5}
\end{figure}

\subsection{Testing the correction}
\label{sec:c5}
\textbf{Code:} \code{experiments/task4\_prior\_shift.py}.

We tested the correction on the two models of Section~\ref{sec:models}, using data where only the
spam rate changes, with $\cfp=10$, $\cfn=1$,
$\ci=0.3$. On synthetic data the models are trained at $\pio=0.5$ and tested on fresh samples at
each deployment rate $\pion$ from $0.05$ to $0.8$. On Spambase (trained at $\pio=0.394$) the test part
is resampled within each class to the new rate, for example $586$ emails of which $29$ are spam at
$\pion=0.05$. Following Section~\ref{sec:gaussian}, the Gaussian model is the shared one with $\log(1+x)$
features on Spambase. The decisions are made on four versions of $\post$: uncorrected, corrected with the true $\pion$, corrected with $\pion$
estimated from the unlabelled test emails (an EM iteration that sets the rate to the average corrected
probability until it stops changing), and
(synthetic only) the true posterior.

\begin{figure}[!htbp]
\centering
\includegraphics[width=0.88\linewidth]{c5_prior_shift.pdf}
\caption{Realized loss against the deployment spam rate $\pion$, without correction, with the known
$\pion$ and with $\pion$ estimated by EM (bands: $\pm1$ standard error over 20 splits). Dashed: the
true posterior. Dotted: the training rate.}
\label{fig:c5}
\end{figure}

\begin{table}[!htbp]
\centering\footnotesize
\input{\resultsdir/c5_loss_spambase.tex}
\caption{Spambase: realized loss after the spam rate changes to $\pion$ (training rate $0.394$),
without correction, corrected with the known $\pion$, and corrected with the EM estimate
(mean $\pm$ standard error, 20 splits).}
\label{tab:c5}
\end{table}

Figure~\ref{fig:c5} and Table~\ref{tab:c5} show the following.
\begin{itemize}
  \item \emph{The correction pays off, and more so the larger the shift.} At $\pion=0.05$ it more than
        halves the loss: on Spambase from $0.095$ to $0.038$ (logistic) and from $0.108$ to $0.051$
        (Gaussian); on synthetic data from $0.154$ and $0.137$ to $0.040$, close to the $0.034$ of the
        true posterior. At the training rate the curves coincide, as they should.
  \item \emph{Where the saving comes from.} Without correction the filter still believes half of the
        synthetic emails are spam, so it inspects $43$--$49\%$ of them; after correction it inspects
        $1.6\%$. The corrected filter is slightly \emph{less} accurate ($96.4\%$ against $98.7\%$ on the
        emails it does not inspect) but three to four times cheaper.
  \item \emph{An unknown rate can be estimated if the model is calibrated.} For logistic regression on
        Spambase the EM estimate is off by $0.007$--$0.013$ on average and the loss is almost the same as
        with the true rate. The Gaussian model's estimate is too low: by $0.04$--$0.07$ on synthetic data,
        and by $0.036$--$0.047$ on Spambase once $\pion\ge0.4$.
\end{itemize}
Again the more accurate filter was not the cheaper one. The last part of the project measured this
trade-off directly, together with the value of inspection.
